\documentclass[11pt,letterpaper]{article}
\usepackage[margin=0.7in,headheight=15pt]{geometry}
\usepackage[T1]{fontenc}\usepackage{lmodern,amsmath,amssymb,fancyhdr,parskip,hyperref}
\hypersetup{colorlinks=true,urlcolor=blue}
\pdfinfoomitdate=1\pdftrailerid{}\pdfsuppressptexinfo=15
\pagestyle{fancy}\fancyhf{}\lhead{Week 81 / Buckets of Fish}\rhead{Facilitator draft}\cfoot{Bellingham Math Circle / Week 81 / facilitator / \thepage}
\setlength{\parskip}{7pt}\setlength{\parindent}{0pt}
\begin{document}
\section*{Mathematical overview}
With finitely many ordered buckets, a move removes exactly one counter from a nonempty bucket and may add any \emph{finite} number only to buckets strictly to its left. Every play terminates, even if total counter count increases. The rightmost bucket never refills and is reduced only finitely often; after its last reduction the lower-bucket game terminates by induction. Finite refills and a fixed finite number of buckets are essential.

There is no uniform finite length bound for the uncapped game from $(0,1)$: remove the right fish and refill the left with $M$ counters, producing a play of exactly $M+1$ moves. Each such play is finite.

Under normal play, exactly the all-even positions are losing. From an all-even position any move produces an odd bucket. From any other position, remove from the rightmost odd bucket and add 0 or 1 to each earlier odd bucket to restore all-even. The same parity argument works when the warm-up permits at most one addition to each earlier bucket. No infinite play and no optimal-duration claim is implied.

Grades 2--3 play and test; Grades 4--5 can explain termination/unbounded length, then parity. Experiments suggest a strategy; the two-direction argument establishes it.
\newpage
\section*{Materials and common launch}
For each pair: print page 2's reusable board, page 3's sixteen sorting cards, and page 6's large-quantity cards when using the upper continuation. Cut the cards in advance; put quantity cards in erasable sleeves and supply a dry-erase pen, or use pencil and eraser. Supply thirty loose counters and a reserve. For two-bucket tasks use only columns 1 and 2 and cover column 3; left/right order is unchanged. Use counters OR a quantity card as the single count in a bucket, never both. Erase and rewrite a count card after each removal or refill; the card's numeral represents that finite many counters. Put the supply outside the board; it is not a fourth bucket. The upper triplet rotates player/player/referee. One anchored adult per participating table is realistic. Cards can stand for a named large \emph{finite} quantity during a thought experiment.

Reading: short rules aloud if needed. Arithmetic: decrement by one, small counts; pairing counters for evenness. Reasoning: maintain left-to-right order and distinguish a chosen large finite quantity from infinitely many counters. No exponent notation is needed by children.

Allow materials play after the usual movement break. Launch in three minutes: demonstrate removing one from the middle, optionally adding only left, then showing the resulting state. Partner points to the bucket selected before removal and checks each addition. The printed worked example makes this recoverable. Never permit rightward or same-bucket refills.

Grades 2--3 first route: 15 minutes capped Problem 1 games with swapped starters, 15 minutes sorting Problem 2 starts, then Problem 3 comparisons if legal moves are fluent. Parity is optional. Grades 4--5 continuation: Problems 4--5 explicitly remove the refill cap; allow 15 minutes trying to make the two-bucket game last and 15 minutes explaining why every chosen play ends. One main investigation is enough for a meeting; uncapped termination is not required of every child. Use fixed small counters for physical play; a quantity card does not require the child to make hundreds of moves. All timing and materials handling untested.
\newpage
\section*{Solutions, hints and optional strategy}
Problem 1: $(0,0,1)$, $(2,1,1)$ and $(1,1,0)$ are first-player wins; $(2,2,2)$ is a second-player win. Problem 2: exactly $(0,0),(0,2),(2,0),(2,2)$ are second-player wins among the sixteen cards; every other start is first-player. Problem 3: $(2,2,2)$ and $(0,2,0)$ are second-player wins; the other four printed starts are first-player. Pairing each bucket's counters makes the even target visible.

Problem 5, termination without notation: if a three-bucket play were infinite, its rightmost bucket would eventually be unchanged because it can only decrease. There would then be an infinite two-bucket suffix, contradicting the same argument at two buckets. One bucket terminates by its decreasing finite count. The statement concerns an entire play, not a known deadline for the final right-bucket action.

Problem 4: exact lengths 10, 30 and 100 come from first refilling the left bucket with 9, 29 and 99. To beat any requested length $L$, use two buckets initially $(0,1)$ and add $L$ counters left on the first move. The resulting game lasts $L+1$ moves. Adding a larger finite supply lengthens a play; it never creates an infinite play. The warm-up cap prevents this particular arbitrary-length construction, so do not silently apply it to the capped rules.

Parity: let $j$ be the rightmost odd bucket. Reduce it by one. All later buckets were already even and cannot change. Every earlier odd bucket is repaired by adding one, and earlier even buckets by adding zero. Thus there is a legal move to all-even. If you start all-even, every opponent move creates at least one odd bucket; reply this way. Termination ensures the last move reaches the empty all-even position and wins.

Useful questions when needed: which bucket can never refill? If we stop touching that bucket, what smaller game remains? Can you pair every bucket's counters without leftovers? These are adult hints, not mandatory child bookkeeping.

Ordinal description for adults: with left-to-right counts $(a_0,\ldots,a_{m-1})$, use $\omega^{m-1}a_{m-1}+\cdots+\omega a_1+a_0$. Every move decreases it lexicographically from the highest altered bucket. This is ordinary ordinal order, not surreal arithmetic.
\newpage
\section*{Encore, provenance and observation record}
Encore: allow removing 1, 2 or 3 from the chosen bucket, and explicitly lift the warm-up's one-per-bucket refill cap. Permit any finite leftward refill (allowing at least 0 through 3 is enough for the strategy). Investigate multiples of four as the new losing target. The proof parallels parity: change the rightmost nonmultiple to a multiple, then adjust all earlier counts upward by at most three. Keep this rule change with adults until the basic game is fluent.

A second return visit can compare fixed game length in a one-bucket start with variable length in two buckets. Children can design starts that look larger but finish sooner; total counter count is not a decreasing potential in the general game.

Source: Joel David Hamkins, ``Buckets of fish!'' (March 4, 2017), \href{https://jdh.hamkins.org/buckets-of-fish/}{Buckets of fish!}. The game, finite-refill termination and parity strategy are credited to that presentation; our boards, examples and task sequence are new teaching adaptations. Do not reproduce the blog's figures or claim this is a newly invented game.

Pedagogy: Givental, Nemirovskaya and Zakharevich, \emph{Math Circle by the Bay}, preface pp. ix--x: manipulatives and varied pace. Our capped-to-uncapped progression and table roles are proposals, not tested findings from the source. Relevant existing library: Week 8 Nim and Week 11 chip-firing, which do not already supply this unbounded finite-length kernel.

Unpiloted. Print US Letter, single-sided, Actual Size. Counter fit, physical handling and procedure rehearsal have not been performed. Record actual rule mistakes, unusually long games, children's ending arguments and whether parity was discovered or supplied. Keep observations distinct from hypotheses.
\end{document}
